\documentclass[../../main.tex]{subfiles}
\begin{document}

{\bf [解]}
新しく
\begin{align}
 A=\displaystyle\int_0^1 f(t)e^{-t}dt \label{eq:1}
\end{align}
とおく．題意の関数方程式は
\begin{align}
 \int_{0}^{x} f(t)dt = e^x -aAe^{2x} \label{eq:2}
\end{align}
である．両辺を$x$で微分して
\begin{align}
f(x)=e^x-2aA\,e^{2x}  \label{eq:3}
\end{align}
だから，\cref{eq:1}に代入して
\begin{align}
A 
 &=\int_0^1(e^t-2aA\cdot e^{2t})e^{-t}dt \\
 &=\int_0^1(1-2aA\cdot e^t)dt \\
 &=\left[t-2aA\cdot e^t\right]_0^1 \\
 &=1-2aA(e-1)
\therefore\ 
A
 &=\frac{1}{1+2a(e-1)} \label{eq:4}
\end{align}
となる．

次に\cref{eq:2}に $x=0$ を代入すると
\begin{align}
 0=1-Aa \quad \therefore\ Aa=1 \label{eq:5}
\end{align}
である．
\cref{eq:4,eq:5}より
\begin{align}
(A,a)=\left(3-2e,\ \frac{1}{3-2e}\right)
\end{align}
であり，求める関数$f(x)$は\cref{eq:3}より
\begin{align*}
f(x)=e^x-2e^{2x}, 
\end{align*}
となる．

\newpage

\end{document}
